% ---------------------------------------------------------------
% TEMPLATE TCC1 - SENAC SANTO AMARO
% ---------------------------------------------------------------
% Trabalho de Conclusão de Curso 1
% Elaborado para apresentação da proposta inicial do TCC
% Estrutura: 4 capítulos (Introdução, Revisão Bibliográfica,
%            Desenvolvimento, Referências Bibliográficas)
% ---------------------------------------------------------------

\documentclass[
    12pt,                % tamanho da fonte
    openright,           % capítulos começam em página ímpar
    oneside,             % para impressão em apenas um lado do papel
    a4paper,             % tamanho do papel
    english,             % idioma adicional
    brazil               % idioma principal
]{abntex2}

% ---------------------------------------------------------------
% PACOTES
% ---------------------------------------------------------------
\usepackage{lmodern}                % Fonte Latin Modern
\usepackage[T1]{fontenc}            % Seleção de códigos de fonte
\usepackage[utf8]{inputenc}         % Codificação do documento
\usepackage{lastpage}               % Usado pela Ficha catalográfica
\usepackage{indentfirst}            % Indenta o primeiro parágrafo
\usepackage{color}                  % Controle de cores
\usepackage{graphicx}               % Inclusão de gráficos
\usepackage{microtype}              % Melhorias de justificação
\usepackage{float}                  % Para posicionamento de figuras
\usepackage{booktabs}               % Tabelas profissionais
\usepackage{multirow}               % Células mescladas em tabelas
\usepackage{longtable}              % Tabelas longas
\usepackage{amsmath,amssymb}        % Símbolos matemáticos
\usepackage{pgfgantt}               % Diagramas de Gantt
\usepackage[brazilian,hyperpageref]{backref}  % Páginas com citações
\usepackage[alf]{abntex2cite}       % Citações ABNT

% ---------------------------------------------------------------
% CONFIGURAÇÕES DO PDF
% ---------------------------------------------------------------
\makeatletter
\hypersetup{
    pdftitle={\@title},
    pdfauthor={\@author},
    pdfsubject={TCC1 - SENAC Santo Amaro},
    pdfcreator={LaTeX with abnTeX2},
    pdfkeywords={tcc}{senac}{santo amaro}{trabalho de conclusão},
    colorlinks=true,        % false: links em caixas; true: links coloridos
    linkcolor=blue,         % cor dos links internos
    citecolor=blue,         % cor dos links para bibliografia
    filecolor=magenta,      % cor dos links para arquivos
    urlcolor=blue,
    bookmarksdepth=4
}
\makeatother

% ---------------------------------------------------------------
% CONFIGURAÇÕES DE APARÊNCIA DO PDF FINAL
% ---------------------------------------------------------------
\definecolor{blue}{RGB}{41,5,195}

% ---------------------------------------------------------------
% INFORMAÇÕES DO DOCUMENTO
% ---------------------------------------------------------------
\titulo{Detecção de vulnerabilidades de memória em C: Abordagem integrando VSA e HOFG}
\autor{Lucas Zacarias Martins Neves}
\local{São Paulo}
\data{2025}
\orientador{Prof. Nome do Orientador}
% \coorientador{Prof. Nome do Coorientador} % Descomente se houver coorientador

\instituicao{%
  SENAC -- Serviço Nacional de Aprendizagem Comercial
  \par
  Unidade Santo Amaro
  \par
  Bacharelado em Ciência da computação
}

\tipotrabalho{Trabalho de Conclusão de Curso I}

\preambulo{Proposta de Trabalho de Conclusão de Curso apresentado ao SENAC Santo Amaro como requisito parcial para aprovação na disciplina de TCC1 do curso de Ciência da computação.}

% ---------------------------------------------------------------
% COMPILA O ÍNDICE
% ---------------------------------------------------------------
\makeindex

% ---------------------------------------------------------------
% INÍCIO DO DOCUMENTO
% ---------------------------------------------------------------
\begin{document}

% Seleciona o idioma do documento
\selectlanguage{brazil}

% Retira espaço extra obsoleto entre as frases
\frenchspacing

% ---------------------------------------------------------------
% ELEMENTOS PRÉ-TEXTUAIS
% ---------------------------------------------------------------
\pretextual

% ---------------------------------------------------------------
% CAPA
% ---------------------------------------------------------------
\imprimircapa

% ---------------------------------------------------------------
% FOLHA DE ROSTO
% ---------------------------------------------------------------
\imprimirfolhaderosto

% ---------------------------------------------------------------
% RESUMO
% ---------------------------------------------------------------
\begin{resumo}

\vspace{\onelineskip}

\noindent
\textbf{Palavras-chave}: Vazamento de memória, HOFG.
\end{resumo}

% ---------------------------------------------------------------
% LISTA DE ILUSTRAÇÕES
% ---------------------------------------------------------------
% \pdfbookmark[0]{\listfigurename}{lof}
% \listoffigures*
% \cleardoublepage

% ---------------------------------------------------------------
% LISTA DE TABELAS
% ---------------------------------------------------------------
% \pdfbookmark[0]{\listtablename}{lot}
% \listoftables*
% \cleardoublepage

% ---------------------------------------------------------------
% SUMÁRIO
% ---------------------------------------------------------------
\pdfbookmark[0]{\contentsname}{toc}
\tableofcontents*
\cleardoublepage

% ---------------------------------------------------------------
% ELEMENTOS TEXTUAIS
% ---------------------------------------------------------------
\textual

% ===============================================================
% CAPÍTULO 1 - INTRODUÇÃO
% ===============================================================
\chapter{Introdução}
\label{cap:introducao}

\section{Objetivos}
\label{sec:Objetivos}

\subsection{Objetivo geral}

O objetivo principal do trabalho é implementar um analisador estático de código para a linguagem C, com foco na detecção de vulnerabilidades de memória, através da integração da representação Heap Object Flow Graph (HOFG) com técnicas de Value-Set Analysis (VSA) para o rastreamento preciso de ponteiros (offset e range).

\subsection{Objetivos específicos}

\begin{enumerate}
    \item Fundamentar teoricamente a integração do Value-Set analysis à estrutura do Heap Object Flow Graph (HOFG), demonstrando como essa abordagem conjunta pode mitigar as limitações de rastreamento de ponteiros e reduzir a taxa de falsos positivos da análise estática.
    \item Implementar motor de análise estática processando a representação intermediária LLVM IR, construindo os algoritmos necessários para a extração e geração do HOFG e sua integração com o VSA.
    \item Avaliar a precisão e a eficiência do analisador implementado utilizando conjuntos de testes que exploram falhas da linguagem C (como vazamentos de memória e usos após desalocação). 
    \item Comparar os resultados obtidos com referências comparativas tendo como foco principal a taxa de falsos positivos e verdadeiro positivos, e comparar o tempo e custo computacional da implementação original do HOFG, para validar o impacto da integração.
\end{enumerate}

% ---------------------------------------------------------------

% ---------------------------------------------------------------
\section{Justificativa}
\label{sec:Hipótese de pesquisa}

A linguagem C tem aplicações no desenvolvimento de sistemas operacionais, sistemas embarcados, drivers, softwares de alta performance, dentre outros. Sendo uma linguagem
que permite o desenvolvedor gerenciar a alocação de itens na memória heap, sem a implementação nativa de elementos tiram tal responsabilidade dos desenvolvedores, encontra como problemas principais os riscos incluindo a má gestão de ponteiros e vazamentos de memória.

Mesmo com a existência de linguagens de programação como Rust, sendo focado em segurança, existem projetos que apresentariam custos excessivos para portabilidade entre as linguagens, fazendo com que seja necessário o uso da linguagem original do desenvolvimento, como o C, neste caso. 
Com isso, para garantir maior segurança de tais projetos, um analisador de código pode ser desenvolvido como uma das soluções, com foco em encontrar potenciais vazamentos e problemas provenientes da má gestão de itens armazenados na memória. 
% ===============================================================
% CAPÍTULO 2 - REVISÃO BIBLIOGRÁFICA
% ===============================================================
\chapter{Revisão Bibliográfica}
\label{cap:revisao}
% ---------------------------------------------------------------
\section{Referencial teórico}
\label{sec:referencial}

% ---------------------------------------------------------------
\section{Metodologia da pesquisa bibliográfica}
\label{sec:metodologia_pesquisa}

Design Science Research
% ---------------------------------------------------------------

\section{Métricas de Avaliação}
\label{sec:metricas_de_avaliacao}

\section{Fundamentos}
\label{sec:fundamentos}

% ---------------------------------------------------------------
\section{Trabalhos relacionados}
\label{sec:trabalhos_relacionados}

O analisador será avaliado a partir do Juliet Test Suite, sendo um conjunto de casos da linguagem C que exploram as suas vulnerabilidades. O resultado obtido será comparado com os resultados de analisadores e algoritmos desenvolvidos que possuem como finalidade o mesmo objetivo, ou seja, de garantir a segurança dos projetos. O principal analisador a ser comparado é o HOFG-Analyser, sendo a aplicação original da representação de programas.


% ===============================================================
% CAPÍTULO 3 - DESENVOLVIMENTO
% ===============================================================
\chapter{Desenvolvimento}
\label{cap:desenvolvimento}
% ---------------------------------------------------------------
\section{Solução proposta}
\label{sec:solucao}

% ---------------------------------------------------------------
\section{Cronograma de trabalho}
\label{sec:cronograma}


% ---------------------------------------------------------------
% ELEMENTOS PÓS-TEXTUAIS
% ---------------------------------------------------------------
\postextual

% ===============================================================
% CAPÍTULO 4 - REFERÊNCIAS BIBLIOGRÁFICAS
% ===============================================================
% As referências bibliográficas são geradas automaticamente a partir
% do arquivo bibliografia.bib, seguindo o padrão ABNT (NBR 6023).
%
% COMO CITAR no texto:
%   \cite{chave}         -> (SOBRENOME, ano) — citação indireta
%   \citeonline{chave}   -> Sobrenome (ano) — citação direta no texto
%
% TIPOS DE REFERÊNCIA no arquivo .bib:
%   @book        -> Livros
%   @article     -> Artigos de periódicos
%   @inproceedings -> Artigos em anais de eventos
%   @mastersthesis -> Dissertações de mestrado
%   @phdthesis   -> Teses de doutorado
%   @misc        -> Sites, vídeos e outros materiais
%   @manual      -> Documentação técnica
%
% Veja o arquivo bibliografia.bib para exemplos de cada tipo.
% ---------------------------------------------------------------
\bibliography{bibliografia}

% ---------------------------------------------------------------
% APÊNDICES (se necessário)
% ---------------------------------------------------------------
% \begin{apendicesenv}
% \partapendices
%
% \chapter{Nome do Apêndice}
% \label{ap:apendice1}
%
% Conteúdo do apêndice (material elaborado pelo próprio autor).
%
% \end{apendicesenv}

% ---------------------------------------------------------------
% ANEXOS (se necessário)
% ---------------------------------------------------------------
% \begin{anexosenv}
% \partanexos
%
% \chapter{Nome do Anexo}
% \label{an:anexo1}
%
% Conteúdo do anexo (material de terceiros).
%
% \end{anexosenv}

\end{document}
